\documentclass[11pt, a4paper]{article}

\usepackage[T1]{fontenc}
\usepackage[utf8]{inputenc}
\usepackage{tgpagella}
\usepackage[spanish, es-noshorthands]{babel}
\usepackage{amsmath, amssymb, bm}
\usepackage{tabularx}
\usepackage{booktabs}
\usepackage[most]{tcolorbox}
\usepackage[margin=2.5cm]{geometry}
\usepackage{ragged2e}
\usepackage{fancyhdr}

\pagestyle{fancy}
\fancyhf{}
\setlength{\headheight}{16pt}
\lhead{\textbf{UCB Sede Tarija} | Ing. Mecatrónica}
\chead{}
\rhead{IMT-121 Circuitos Electrónicos I | Registro Lab B}
\lfoot{Prof: M.Sc. Ing. Bernardo Quiroga Turdera}
\rfoot{Página \thepage}
\renewcommand{\headrulewidth}{0.4pt}

\definecolor{ucbblue}{RGB}{0, 51, 102}

\begin{document}

\begin{tcolorbox}[colback=ucbblue!5!white, colframe=ucbblue, title=\textbf{Hoja de Registro de Datos: Experiencia Integrada B}]
    \centering \textbf{Registro de Barrido de Cargas, Potencia y Discrepancias}
\end{tcolorbox}

\textbf{Equipo:} \hrulefill

\vspace{0.2cm}
\textbf{Valores Base del Equivalente:} $V_{th} = $ \rule{2cm}{0.4pt} V \quad | \quad $R_{th} = $ \rule{2cm}{0.4pt} k$\Omega$

\section*{Tabla de Barrido de Potencia}
\textit{Nota: Registre si $I_L$ es medida directamente o calculada de manera indirecta ($V_L/R_L$).}

\vspace{0.3cm}
\renewcommand{\arraystretch}{2.2}
\noindent
% Ajustamos las proporciones de las columnas X. Suma total de \hsize = 7.
% 6 columnas * 0.85 = 5.1. La ultima columna = 1.9. Total = 7.
\begin{tabularx}{\textwidth}{*{6}{>{\centering\arraybackslash\hsize=0.85\hsize}X} >{\RaggedRight\arraybackslash\hsize=1.9\hsize}X}
\toprule
\textbf{$\bm{R_L}$ med.} \newline ($\Omega$) & 
\textbf{$\bm{V_L}$ SPICE} \newline (V) & 
\textbf{$\bm{V_L}$ med.} \newline (V) & 
\textbf{$\bm{I_L}$ deriv.} \newline (mA) & 
\textbf{$\bm{P_L}$ SPICE} \newline (mW) & 
\textbf{$\bm{P_L}$ med.} \newline (mW) & 
\textbf{Observación} \newline (Error \%) \\ 
\midrule
 & & & & & & \\ \midrule
 & & & & & & \\ \midrule
 & & & & & & \\ \midrule
 & & & & & & \\ \midrule
 & & & & & & \\ \midrule
 & & & & & & \\ \midrule
 & & & & & & \\ 
\bottomrule
\end{tabularx}

\vspace{0.3cm}
\section*{Análisis de Resultados (Efecto de Carga y Máximo)}
\begin{enumerate}
    \item Resistencia experimental donde se obtuvo el mayor valor de potencia ($P_{L,\text{max}}$): 
    \par\vspace{0.2cm}\noindent\hrulefill
    
    \item Desviación respecto al valor teórico predicho ($\Delta R = R_{L,\text{max,exp}} - R_{th}$): 
    \par\vspace{0.2cm}\noindent\hrulefill
    
    \item \textbf{Efecto de Carga:} Observe sus datos de $V_L$. ¿Qué ocurre con la tensión que recibe la carga a medida que inserta resistencias ($R_L$) cada vez más grandes? Justifique brevemente mediante el concepto de divisor de tensión con $R_{th}$.
    \vspace{2cm}
    
    \item A partir de sus observaciones, formule \textbf{una hipótesis técnica} que justifique las discrepancias entre SPICE y la medición de potencia máxima obtenida. (Evite justificaciones como "error humano").
    \vspace{2cm}
\end{enumerate}

\end{document}
